% !TeX program = pdflatex
%==============================================================================
%  Verificação Suplementar (VS) --- prova final, aulas 01 a 10
%
%  O gabarito sai de "Verificacao suplementar - gabarito.tex", que apenas liga a
%  opção [gabarito] e inclui este arquivo --- fonte única, dois PDFs.
%
%  QUATRO QUESTÕES, DE 3,0 PONTOS: duas sorteadas da "Lista de revisao 01-06" e
%  duas da "Lista de revisao 07-10", uma aula por questão dentro de cada lista.
%  Ficaram fora do sorteio os exercícios que já tinham caído na Avaliação Teórica
%  01 (Ex. 5, 7 e 11 da 01-06) e na 02 (Ex. 3, 6 e 9 da 07-10). A Verificação de
%  Reposição (VR) e a VS foram sorteadas juntas, a VR primeiro, e não têm
%  exercício em comum. O sorteio é reproduzível:
%
%      import numpy as np
%      rng = np.random.default_rng(20261116)                 # a data, aaaammdd
%      # exercícios que ainda não caíram em prova, por aula
%      lista_01_06 = {1: [1, 2], 2: [3, 4], 3: [6], 4: [8], 5: [9, 10]}
%      lista_07_10 = {7: [1, 2], 8: [4], 9: [5], 10: [7, 8]}
%
%      def sorteia(lista):
%          aulas = sorted(rng.choice(sorted(lista), size=2, replace=False).tolist())
%          escolhidos = []
%          for a in aulas:
%              e = int(rng.choice(lista[a]))
%              escolhidos.append(e)
%              lista[a].remove(e)
%              if not lista[a]:
%                  del lista[a]
%          return escolhidos
%
%      vr = sorteia(lista_01_06) + sorteia(lista_07_10)
%      vs = sorteia(lista_01_06) + sorteia(lista_07_10)
%      # -> vr = [1, 4, 2, 5] (aulas 01, 02, 07 e 09)
%      #    vs = [2, 3, 1, 7] (aulas 01, 02, 07 e 10)
%
%  Resultado, para a VS:
%      Aula 01 -> Ex. 2 da 01-06  ex-02-dois-riscos (risco condicional e risco
%                                 esperado)
%      Aula 02 -> Ex. 3 da 01-06  ex-03-ortonormal (Ridge e Lasso no caso
%                                 ortonormal)
%      Aula 07 -> Ex. 1 da 07-10  ex-01-logistica-sem-gaussianas (a logística
%                                 sem as hipóteses do LDA)
%      Aula 10 -> Ex. 7 da 07-10  ex-07-voto-vizinhos (o voto de k vizinhos)
%
%  O sorteio rodou uma vez, em 16/11/2026, e vale. Como a VS só podia usar o que
%  a VR deixou, as duas cobrem as aulas 01 e 02 na parte de regressão.
%
%  NÃO HÁ TEXTO DE EXERCÍCIO NESTE ARQUIVO. Os enunciados e as soluções vêm por
%  \input dos MESMOS arquivos que as listas de revisão usam, em exercicios/ e em
%  exercicios-07-10/. Os pontos por item escritos neles somam 4,0, o peso de uma
%  questão nas provas de três questões; aqui, o \pontositens antes de cada
%  \input dá os valores desta prova, que somam 3,0.
%
%  A PROVA NÃO REVELA DE QUE AULA VEM CADA QUESTÃO: o \questao não leva marcador
%  de aula, os títulos ficam vazios, a pedido do Gabriel, e a
%  caixa de instruções traz a notação de regressão e a de classificação juntas.
%==============================================================================
\providecommand{\gabaritoopt}{}
\documentclass[11pt,a4paper]{article}
\usepackage[\gabaritoopt]{../recursos/latex/estilo-avaliacao}

\provatrue
\renewcommand{\tituloprova}[1]{Verificação Suplementar (#1)}

\begin{document}
\cabecalhoprova[{salvo indicação em contrário, $(\X,Y)$ é um vetor aleatório com
            $\X\in\R^p$, e $p$ é o número de covariáveis;
            $\Dados=\{(\X_i,Y_i)\}_{i=1}^n$ é a amostra de treino, i.i.d. Em
            regressão, $r(\x)=\E[Y\mid\X=\x]$ é a função de regressão, e a perda é a
            quadrática. Em classificação, $Y$ assume valores num conjunto finito
            $\mathcal{C}$ de classes; no caso binário, $\mathcal{C}=\{0,1\}$, a
            classe $Y=1$ é a positiva e $\eta(\x)=\Prob(Y=1\mid\X=\x)$, e a perda é
            a 0--1, com risco $\risco(g)=\Prob\big(Y\ne g(\X)\big)$.}]{VS}{01 a 10}{Regressão e Classificação}{2 horas}{12,0 pontos}

%==============================================================================
\questao{}{3,0 pontos}
%==============================================================================
\pontositens{1,0;0,9;1,1}
% Aula 01 --- risco condicional x risco esperado, e a barra de erro do teste.
A palavra "risco" apareceu durante o nosso curso com dois sentidos, e a diferença entre eles
não é preciosismo: ela decide qual pergunta o seu número está respondendo.

Quando $g$ é tratada como uma função \emph{fixa}, a esperança em
$\E[(Y-g(\X))^2]$ percorre só a observação nova $(\X,Y)$. Chamamos isso de
\textbf{risco condicional}. Quando $g=\ghat$ é tratada como aleatória --- e ela é,
porque foi construída a partir de $\Dados$ ---, a esperança percorre também as
amostras de treino possíveis, e temos o \textbf{risco esperado}.
\begin{itens}
  \item \pts{1,3} Cada um dos dois oferece uma garantia diferente. Descreva as
        duas, e diga qual interessa a quem já ajustou um modelo e vai colocá-lo em
        produção, e qual interessa a quem ainda está escolhendo entre diversos métodos.
  \item \pts{1,2} Você separou um conjunto de teste de $m$ observações que não
        participou de nenhum treinamento, e calculou $\riscohat(\ghat)=\frac1m\sum_{k=1}^m W_k$,
        com $W_k=(Y_k-\ghat(\X_k))^2$. Qual dos dois riscos esse número estima?
        Para responder, pergunte-se o que está fixo e o que varia no momento em
        que você faz a conta.
  \item \pts{1,5} Um conjunto de teste permite algo que a validação cruzada não
        permite: uma barra de erro honesta. Explique por que os $W_k$ são i.i.d.\ e
        escreva o intervalo de confiança aproximado de 95\% que o Teorema Central
        do Limite fornece.
\end{itens}

\begin{solucao}
\textbf{(a)} O \textbf{risco condicional} é uma garantia sobre \emph{aquela
função específica}: se eu usar este $\ghat$, já ajustado, em observações novas, meu
erro quadrático médio esperado é tanto.

O \textbf{risco esperado} é uma garantia sobre o \emph{procedimento}: se eu
sortear uma amostra de treino, rodar este método e usar o que sair, meu erro
esperado é tanto. A esperança percorre também os treinos possíveis.

Quem vai pôr o modelo em produção quer o \textbf{condicional}. O modelo dele é um
só, já existe, não vai ser reajustado --- saber como se comportam os modelos que
aquele método \emph{costuma} produzir não paga as contas.

Quem está comparando métodos quer o \textbf{esperado}. A pergunta ali é "qual
deles tende a produzir bons modelos", e uma rodada de sorte com um método ruim não
deveria decidir a escolha.

\smallskip
\textbf{(b)} Estima o \textbf{risco condicional}.

No momento da conta, $\ghat$ já foi ajustado e não muda: é uma função fixa. O que
varia de uma parcela para outra é apenas o par $(\X_k,Y_k)$ sorteado para o
teste. A aleatoriedade do treino --- que outra amostra teria dado outro $\ghat$
--- não entra na conta em momento nenhum, porque o treino aconteceu antes e não é
repetido.

Dito de outro modo: para estimar o risco \emph{esperado} você teria de repetir o
processo inteiro, sorteando vários treinos, e não é o que um único conjunto de
teste faz.

\smallskip
\textbf{(c)} Condicionalmente a $\ghat$ --- que é a situação, por (b) --- cada
$W_k=(Y_k-\ghat(\X_k))^2$ é função de um único par $(\X_k,Y_k)$. Esses pares são
i.i.d.\ entre si e independentes do treino, e a mesma função $\ghat$ é aplicada a
todos. Logo os $W_k$ são i.i.d., com média $\risco(\ghat)$.

Como $\riscohat$ é média de i.i.d., o Teorema Central do Limite se aplica e dá
\[
  \riscohat(\ghat)\;\pm\;2\sqrt{\widehat\sigma^2/m},
  \qquad
  \widehat\sigma^2=\frac1m\sum_{k=1}^m(W_k-\overline W)^2 .
\]

Há uma hipótese fazendo todo o trabalho aqui: o teste não participou de nada. Se ele
tivesse sido usado para escolher qualquer coisa, os $W_k$ deixariam de ser
independentes do $\ghat$ e o intervalo perderia o sentido.
\end{solucao}


\begin{criterio}
\textbf{(a)} 0,3 pela garantia do risco condicional, sobre a função ajustada que se tem
nas mãos; 0,3 pela do risco esperado, sobre o procedimento, em média sobre as amostras de
treino; 0,4 por dizer a quem cada um interessa, 0,2 cada: o condicional a quem põe em
produção um modelo já ajustado, o esperado a quem escolhe entre métodos.

\textbf{(b)} 0,3 por identificar o que está fixo no momento da conta (o treino e,
portanto, $\ghat$) e o que varia (os pares de teste); 0,3 pela esperança condicional a
$\Dados$, igual a $\risco(\ghat)$, ou pela lei dos grandes números condicional; 0,3 pela
resposta, o risco condicional. A observação de que o erro de teste também é não viesado
para o risco esperado, mas não converge para ele, é bem-vinda e não é exigida; responder
``o risco esperado'' só com o argumento da não viesadez vale 0,3 do item.

\textbf{(c)} 0,5 por justificar que os $W_k$ são i.i.d.\ condicionalmente a $\Dados$: com
o treino fixo, $\ghat$ é uma função fixa, e $W_k$ é a mesma função aplicada a pares i.i.d.\
e independentes do treino. Dizer apenas que os dados de teste são i.i.d.\ vale 0,2 desses
0,5. 0,6 pelo intervalo: 0,3 pelo Teorema Central do Limite aplicado à média dos $W_k$,
com a variância estimada pela variância amostral, e 0,3 pela forma
$\riscohat(\ghat)\pm2\sqrt{\widehat\sigma^2/m}$ (aceitar $1{,}96$ no lugar de $2$). A
comparação com a validação cruzada não é exigida.
\end{criterio}

%==============================================================================
\questao{}{3,0 pontos}
%==============================================================================
\pontositens{0,8;0,5;1,1;0,6}
% Aula 02 --- derivar as fórmulas do Ridge e do Lasso no caso ortonormal.
As notas da aula de Regressão Linear apresentam duas fórmulas prontas para o caso ortonormal,
$\mathbb{X}^\top\mathbb{X}=\mathbb{I}$, e não as derivam. São elas que explicam a
frase "o Ridge encolhe e o Lasso zera", e o exercício é preencher essa lacuna.

Use as definições da aula, \textbf{sem} fator $\tfrac12$ no RSS:
\[
  \widehat{\bbeta}^{\,\text{ridge}}
   = \argmin_{\bbeta}\Big\{\norm{\bm y-\mathbb{X}\bbeta}^2+\lambda\sum_j\beta_j^2\Big\},
  \qquad
  \widehat{\bbeta}^{\,\text{lasso}}
   = \argmin_{\bbeta}\Big\{\norm{\bm y-\mathbb{X}\bbeta}^2+\lambda\sum_j\abs{\beta_j}\Big\}.
\]
\begin{itens}
  \item \pts{1,1} Comece pelo passo que torna tudo possível. Supondo
        $\mathbb{X}^\top\mathbb{X}=\mathbb{I}$, mostre que o estimador de mínimos
        quadrados é $\widehat\bbeta=\mathbb{X}^\top\bm y$ e que
        \[
          \norm{\bm y-\mathbb{X}\bbeta}^2
            = \norm{\bm y}^2-\norm{\widehat\bbeta}^2
              + \sum_{j=1}^{p}\big(\beta_j-\widehat\beta_j\big)^2 .
        \]
        Explique por que essa identidade faz o problema \textbf{se separar} em $p$
        problemas de uma variável --- e por que é justamente isso que falha quando
        as colunas não são ortonormais.
  \item \pts{0,7} Com a separação em mãos, derive
        $\widehat\beta_j^{\,\text{ridge}}=\widehat\beta_j/(1+\lambda)$.
  \item \pts{1,4} Agora o Lasso. A dificuldade nova é que $\abs{\beta_j}$ não é
        derivável em zero, então não dá para igualar a derivada a zero e pronto:
        trate separadamente as regiões $\beta_j>0$, $\beta_j<0$ e o ponto
        $\beta_j=0$. Derive
        $\widehat\beta_j^{\,\text{lasso}}
          =\operatorname{sinal}(\widehat\beta_j)\big(\abs{\widehat\beta_j}-\lambda/2\big)_+$
        e diga qual é a condição exata que leva um coeficiente a zero.
  \item \pts{0,8} O limiar deu $\lambda/2$, e não $\lambda$. De onde veio esse
        $2$? Diga o que mudaria nas duas fórmulas se a aula tivesse definido o RSS
        com um $\tfrac12$ na frente, e que lição prática isso deixa.
\end{itens}

\begin{solucao}
\textbf{(a)} Com $\mathbb{X}^\top\mathbb{X}=\mathbb{I}$, as equações normais dão
diretamente
$\widehat\bbeta=(\mathbb{X}^\top\mathbb{X})^{-1}\mathbb{X}^\top\bm y
=\mathbb{X}^\top\bm y$.

Agora expanda o RSS, usando essa igualdade e a ortonormalidade:
\[
  \norm{\bm y-\mathbb{X}\bbeta}^2
   = \bm y^\top\bm y - 2\,\bbeta^\top\underbrace{\mathbb{X}^\top\bm y}_{=\;\widehat\bbeta}
     + \bbeta^\top\underbrace{\mathbb{X}^\top\mathbb{X}}_{=\;\mathbb{I}}\bbeta
   = \norm{\bm y}^2 - 2\,\bbeta^\top\widehat\bbeta + \norm{\bbeta}^2 .
\]
Os dois últimos termos pedem um completar quadrados:
\[
  -2\,\bbeta^\top\widehat\bbeta+\norm{\bbeta}^2
   = \norm{\bbeta-\widehat\bbeta}^2-\norm{\widehat\bbeta}^2 ,
\]
o que dá a identidade do enunciado.

O que ela tem de especial: o único termo que depende de $\bbeta$ é
$\sum_j(\beta_j-\widehat\beta_j)^2$, e nessa soma a $j$-ésima parcela envolve
\emph{apenas} $\beta_j$. Como as penalidades $\sum_j\beta_j^2$ e
$\sum_j\abs{\beta_j}$ também são somas coordenada a coordenada, o problema inteiro
vira $p$ problemas escalares que não conversam entre si --- e problema escalar a
gente resolve derivando.

Quando as colunas \emph{não} são ortonormais, $\mathbb{X}^\top\mathbb{X}$ tem
termos fora da diagonal, aparecem produtos $\beta_j\beta_\ell$ com $j\ne\ell$, e a
separação morre. É exatamente por isso que o Lasso não tem solução fechada no caso
geral: a ortonormalidade é uma hipótese muito forte, não uma simplificação
inocente.

\smallskip
\textbf{(b)} O problema escalar é minimizar
$h(\beta)=(\beta-\widehat\beta_j)^2+\lambda\beta^2$, que é derivável em toda a
reta:
\[
  h'(\beta)=2(\beta-\widehat\beta_j)+2\lambda\beta=0
  \;\Longrightarrow\;
  \beta(1+\lambda)=\widehat\beta_j
  \;\Longrightarrow\;
  \beta=\frac{\widehat\beta_j}{1+\lambda}.
\]
Como $h''=2(1+\lambda)>0$, é mínimo.

Leia o resultado: o Ridge \textbf{multiplica} por um fator $1/(1+\lambda)<1$. Isso
encolhe proporcionalmente --- quem era grande perde mais em valor absoluto --- e
nunca zera, porque o fator só alcança zero no limite $\lambda\to\infty$.

\smallskip
\textbf{(c)} Agora $h(\beta)=(\beta-\widehat\beta_j)^2+\lambda\abs{\beta}$. Analise
por regiões.

\emph{Onde $\beta>0$}, temos $\abs{\beta}=\beta$ e
$h'(\beta)=2(\beta-\widehat\beta_j)+\lambda$, que zera em
$\beta=\widehat\beta_j-\lambda/2$. Só que essa raiz precisa cair na região em que
a conta vale, isto é, precisa ser positiva --- o que exige
$\widehat\beta_j>\lambda/2$.

\emph{Onde $\beta<0$}, temos $\abs{\beta}=-\beta$ e
$h'(\beta)=2(\beta-\widehat\beta_j)-\lambda$, que zera em
$\beta=\widehat\beta_j+\lambda/2$, admissível se $\widehat\beta_j<-\lambda/2$.

\emph{E se $\abs{\widehat\beta_j}\le\lambda/2$?} Aí nenhuma das duas raízes cai na
sua própria região: a candidata positiva sai negativa e vice-versa. Como $h$ é
convexa e contínua, o mínimo existe, e só pode estar no único ponto que sobrou ---
$\beta=0$, justamente onde a função não é derivável. (Em linguagem de
subgradiente: $0$ é mínimo quando $0\in -2\widehat\beta_j+\lambda[-1,1]$, que é
exatamente a condição $\abs{\widehat\beta_j}\le\lambda/2$.)

Juntando os três casos,
\[
  \widehat\beta_j^{\,\text{lasso}}
   =\operatorname{sinal}(\widehat\beta_j)\big(\abs{\widehat\beta_j}-\lambda/2\big)_+ ,
\]
e a condição exata para zerar é $\boxed{\abs{\widehat\beta_j}\le\lambda/2}$.

Compare com (b) e a diferença salta: o Lasso \textbf{subtrai} uma constante em vez
de multiplicar por um fator. Todo mundo encolhe pela mesma quantidade, $\lambda/2$,
e quem era menor que isso em módulo não sobrevive ao deslocamento. É o
\emph{soft thresholding}, e é a razão algébrica de o Lasso fazer seleção de
variáveis.

\smallskip
\textbf{(d)} O $2$ vem da derivada do quadrado. Derivando
$(\beta-\widehat\beta_j)^2$ aparece um fator $2$; derivando $\lambda\abs{\beta}$
aparece só $\pm\lambda$. Ao isolar $\beta$ dividimos tudo por $2$, e o $\lambda$
vira $\lambda/2$.

Com um $\tfrac12$ na frente do RSS, e só dele, o fator $2$ some da derivada do
quadrado, e a da penalidade não muda. Refazendo as contas, o Lasso fica com
$(\beta-\widehat\beta_j)\pm\lambda=0$, e o limiar vira $\lambda$; o Ridge fica com
$(\beta-\widehat\beta_j)+2\lambda\beta=0$, e a fórmula vira
$\widehat\beta_j/(1+2\lambda)$.

As duas mudanças são, na verdade, uma só. Multiplicar o objetivo inteiro por $2$
não altera o minimizador, e transforma o problema novo no da aula com $2\lambda$ no
lugar de $\lambda$: basta fazer essa troca nas duas fórmulas. Repare que pôr o
$\tfrac12$ também na penalidade não mudaria nada --- seria multiplicar o objetivo
todo por uma constante.

A lição prática é que \textbf{o limiar depende da convenção}, não da matemática.
Antes de comparar um $\lambda$ do seu caderno com um $\lambda$ de um artigo ou de
uma biblioteca, confira qual objetivo cada um está minimizando. O \texttt{Lasso} do
\texttt{scikit-learn} é um bom exemplo: ele minimiza
$\tfrac{1}{2n}\norm{\bm y-\mathbb{X}\bbeta}^2+\alpha\norm{\bbeta}_1$, e o seu
$\alpha$ corresponde a $\lambda=2n\alpha$ na convenção da aula --- o $n$ vem da
média, e o $2$ é exatamente o deste item.
\end{solucao}


\begin{criterio}
\textbf{(a)} 0,2 por $\widehat\bbeta=\mathbb{X}^\top\bm y$, pelas equações normais; 0,3 pela
identidade, por Pitágoras, com a ortogonalidade do resíduo, ou expandindo os quadrados;
0,3 pela separação em $p$ problemas de uma variável, com a justificativa de que uma soma
de funções de coordenadas distintas se minimiza coordenada a coordenada, e por apontar os
termos cruzados que aparecem sem ortonormalidade.

\textbf{(b)} 0,3 pela derivada igualada a zero; 0,2 pela justificativa de que o ponto
crítico é o mínimo, porque a função é estritamente convexa.

\textbf{(c)} 0,3 pelas derivadas fora do zero, nas duas semirretas; 0,6 pelos três casos,
0,2 cada ($\widehat\beta_j>\lambda/2$, $\widehat\beta_j<-\lambda/2$ e
$\abs{\widehat\beta_j}\le\lambda/2$), com o argumento de crescimento e decrescimento dos
dois lados do zero; 0,2 pela fórmula única e pela condição exata de zerar,
$\abs{\widehat\beta_j}\le\lambda/2$. Igualar a derivada a zero sem tratar o ponto
$\beta_j=0$ perde os 0,2 do terceiro caso. Aceitar a solução por subgradientes.

\textbf{(d)} 0,2 pela origem do $2$, a derivada do quadrado; 0,2 pelas duas fórmulas com o
$\tfrac12$ no RSS, $\widehat\beta_j/(1+2\lambda)$ e o limiar $\lambda$, porque multiplicar
o objetivo por $2$ não muda o minimizador; 0,2 pela lição: o valor numérico de $\lambda$
depende da convenção, e é preciso conferir o objetivo de cada fonte antes de comparar. A
correspondência $\lambda=2n\alpha$ do \texttt{scikit-learn} não é exigida.
\end{criterio}

%==============================================================================
\questao{}{3,0 pontos}
%==============================================================================
\pontositens{1,7;1,3}
% Aula 07 --- sob o LDA o log-odds é linear; a recíproca é falsa.
Renata é atuária numa seguradora e quer classificar os segurados de um ramo em dois
perfis, $Y=1$ (risco alto) e $Y=0$ (risco baixo), a partir de uma única covariável: o
tempo $X$, em anos, entre a contratação da apólice e o primeiro sinistro. Ela hesita
entre a análise discriminante linear (LDA) e a regressão logística, dois métodos que
chegam a uma fronteira linear por caminhos diferentes. Ao longo do exercício,
$\sigma(u)=e^u/(1+e^u)$ é a função logística, $\log$ é o logaritmo natural e
$\pi_c=\Prob(Y=c)$.
\begin{itens}
  \item \pts{2,2} Comece pelo caso geral, com $p$ covariáveis: suponha que $\X\in\R^p$ e
        que, como no LDA, $\X\mid Y=c\sim N(\bm\mu_c,\bm\Sigma)$ para $c\in\{0,1\}$, com
        $\bm\Sigma$ simétrica e positiva definida. Mostre que
        $\Prob(Y=1\mid\X=\x)=\sigma(\beta_0+\bbeta^\top\x)$ e escreva $\beta_0$ e
        $\bbeta$ em função de $\pi_0$, $\pi_1$, $\bm\mu_0$, $\bm\mu_1$ e $\bm\Sigma$.
        Conclua que, sob as hipóteses do LDA, o modelo da regressão logística vale
        exatamente. (Escreva $\Prob(Y=1\mid\X=\x)$ pelo Teorema de Bayes, com $f_c$ a
        densidade de $\X$ dado $Y=c$, divida o numerador e o denominador por
        $\pi_1f_1(\x)$ e ponha o resultado na forma $\sigma\big(u(\x)\big)$. Depois,
        substitua as densidades normais em $u(\x)$, expanda as duas formas quadráticas e
        observe o que se cancela porque a covariância é a mesma nas duas classes.)
  \item \pts{1,8} Os tempos de Renata não se parecem com uma normal: são positivos e
        muito assimétricos. Suponha agora que $X\mid Y=c$ tenha distribuição
        exponencial de taxa $\lambda_c$, com densidade $\lambda_c e^{-\lambda_c x}$ para
        $x>0$, com $\lambda_0=1$, $\lambda_1=3$ e $\pi_0=\pi_1=1/2$: em média, o
        primeiro sinistro do perfil de risco alto chega três vezes mais cedo. Mostre que
        $\Prob(Y=1\mid X=x)=\sigma(\log3-2x)$ para $x>0$. O modelo da regressão
        logística vale nesta população, isto é, existem constantes $\beta_0$ e $\beta_1$
        com $\Prob(Y=1\mid X=x)=\sigma(\beta_0+\beta_1x)$ para todo $x>0$? As hipóteses do
        LDA, de que $X$ é normal em cada classe e tem a mesma variância nas duas, valem?
\end{itens}

\begin{solucao}
Em todo o exercício, $0<\pi_1<1$, $\pi_0=1-\pi_1$, e as densidades $f_0$ e $f_1$ de $\X$
dado $Y=0$ e dado $Y=1$ são positivas em todo o domínio considerado: $\R^p$ no item (a), em
que são normais, e $(0,\infty)$ no item (b), em que são exponenciais. Essas hipóteses garantem
que as divisões e os logaritmos abaixo estão bem definidos.

Os dois itens dependem de uma mesma identidade, que convém isolar. Pelo Teorema de Bayes,
\[
  \Prob(Y=1\mid\X=\x)=\frac{\pi_1f_1(\x)}{\pi_0f_0(\x)+\pi_1f_1(\x)} .
\]
Dividindo o numerador e o denominador por $\pi_1f_1(\x)>0$,
\[
  \Prob(Y=1\mid\X=\x)=\frac{1}{1+e^{-u(\x)}}=\sigma\big(u(\x)\big),
  \qquad
  u(\x)=\log\frac{\pi_1f_1(\x)}{\pi_0f_0(\x)},
\]
em que a última igualdade vem de $\sigma(u)=e^u/(1+e^u)=1/(1+e^{-u})$, multiplicando o
numerador e o denominador por $e^{-u}$. A função $\sigma$ é uma bijeção estritamente crescente
de $\R$ em $(0,1)$; por isso $u(\x)$ fica determinado por $\Prob(Y=1\mid\X=\x)$, e é o
\emph{log-odds}, $\log\big[\Prob(Y=1\mid\X=\x)/\Prob(Y=0\mid\X=\x)\big]$. Segue que o modelo
da regressão logística vale, isto é, existem $\beta_0$ e $\bbeta$ com
$\Prob(Y=1\mid\X=\x)=\sigma(\beta_0+\bbeta^\top\x)$ para todo $\x$, se e somente se
$u(\x)=\beta_0+\bbeta^\top\x$ para todo $\x$. Em cada item, basta então decidir se o
\emph{log-odds} é uma função afim de $\x$.

\smallskip
\textbf{(a)} A densidade de $N(\bm\mu_c,\bm\Sigma)$ é
\[
  f_c(\x)=(2\pi)^{-p/2}(\det\bm\Sigma)^{-1/2}
          \exp\Big(-\tfrac12(\x-\bm\mu_c)^\top\bm\Sigma^{-1}(\x-\bm\mu_c)\Big),
\]
positiva em todo $\R^p$. Como a covariância é a mesma nas duas classes, a constante
$(2\pi)^{-p/2}(\det\bm\Sigma)^{-1/2}$ é a mesma em $f_0$ e em $f_1$, e se cancela no
quociente:
\[
  u(\x)=\log\frac{\pi_1}{\pi_0}-\tfrac12\,Q_1(\x)+\tfrac12\,Q_0(\x),
  \qquad
  Q_c(\x)=(\x-\bm\mu_c)^\top\bm\Sigma^{-1}(\x-\bm\mu_c).
\]
A matriz $\bm\Sigma^{-1}$ é simétrica, por ser a inversa de uma matriz simétrica, e por isso
$\x^\top\bm\Sigma^{-1}\bm\mu=\bm\mu^\top\bm\Sigma^{-1}\x$ para todo $\bm\mu$. Expandindo,
\[
  Q_c(\x)=\x^\top\bm\Sigma^{-1}\x-2\,\bm\mu_c^\top\bm\Sigma^{-1}\x
          +\bm\mu_c^\top\bm\Sigma^{-1}\bm\mu_c .
\]
O termo quadrático $\x^\top\bm\Sigma^{-1}\x$ é o mesmo em $Q_0$ e em $Q_1$ e desaparece da
diferença $Q_0-Q_1$. Sobra
\[
  u(\x)=\log\frac{\pi_1}{\pi_0}+(\bm\mu_1-\bm\mu_0)^\top\bm\Sigma^{-1}\x
        -\tfrac12\big(\bm\mu_1^\top\bm\Sigma^{-1}\bm\mu_1-\bm\mu_0^\top\bm\Sigma^{-1}\bm\mu_0\big)
  =\beta_0+\bbeta^\top\x,
\]
com
\[
  \bbeta=\bm\Sigma^{-1}(\bm\mu_1-\bm\mu_0),
  \qquad
  \beta_0=\log\frac{\pi_1}{\pi_0}
     -\tfrac12\big(\bm\mu_1^\top\bm\Sigma^{-1}\bm\mu_1-\bm\mu_0^\top\bm\Sigma^{-1}\bm\mu_0\big),
\]
em que se usou de novo a simetria de $\bm\Sigma^{-1}$, para escrever
$(\bm\mu_1-\bm\mu_0)^\top\bm\Sigma^{-1}\x=\big(\bm\Sigma^{-1}(\bm\mu_1-\bm\mu_0)\big)^\top\x$.
Pela identidade do início, $\Prob(Y=1\mid\X=\x)=\sigma(\beta_0+\bbeta^\top\x)$ para todo
$\x\in\R^p$: sob as hipóteses do LDA, o modelo da regressão logística vale exatamente, com
esses coeficientes.

A conta mostra onde entrou a hipótese de covariância comum. Ela foi usada duas vezes, para
cancelar as constantes de normalização e para cancelar o termo quadrático, e só o segundo uso
é indispensável. Com covariâncias $\bm\Sigma_0\ne\bm\Sigma_1$, as constantes deixariam em
$u(\x)$ a parcela $\tfrac12\log(\det\bm\Sigma_0/\det\bm\Sigma_1)$, que não depende de $\x$ e
iria para o intercepto. Já o termo $-\tfrac12\,\x^\top A\,\x$, com
$A=\bm\Sigma_1^{-1}-\bm\Sigma_0^{-1}$, sobreviveria. A matriz $A$ é simétrica e não nula,
porque inversas iguais viriam de matrizes iguais, e uma matriz simétrica com $\x^\top A\x=0$
para todo $\x$ é nula: com $\x=\bm e_i$ obtém-se $A_{ii}=0$, e com $\x=\bm e_i+\bm e_j$,
$2A_{ij}=0$. O \emph{log-odds} seria quadrático em $\x$, como no QDA.

Quando as hipóteses do LDA valem, as duas abordagens que Renata compara descrevem, portanto, a
mesma família de funções. O LDA estima $\pi_c$, $\bm\mu_c$ e $\bm\Sigma$ e substitui esses
valores nas fórmulas de $\beta_0$ e $\bbeta$; a regressão logística estima $\beta_0$ e $\bbeta$
diretamente.

\smallskip
\textbf{(b)} Para $x>0$, com $f_c(x)=\lambda_ce^{-\lambda_cx}>0$,
\[
  u(x)=\log\frac{\pi_1\lambda_1e^{-\lambda_1x}}{\pi_0\lambda_0e^{-\lambda_0x}}
      =\log\frac{\pi_1}{\pi_0}+\log\frac{\lambda_1}{\lambda_0}-(\lambda_1-\lambda_0)\,x
      =0+\log3-2x ,
\]
e, pela identidade do início, $\Prob(Y=1\mid X=x)=\sigma(\log3-2x)$ para todo $x>0$.

O modelo da regressão logística vale nesta população: o \emph{log-odds} é afim em $x$, com
$\beta_0=\log3\approx1{,}0986$ e $\beta_1=-2$. O sinal de $\beta_1$ diz que, quanto mais tarde
vem o primeiro sinistro, menos provável é o perfil de risco alto.

As hipóteses do LDA, por outro lado, não valem, por dois motivos independentes. Primeiro, $X$
não é normal em nenhuma das classes: $\Prob(X\le0\mid Y=c)=0$, enquanto uma normal dá
probabilidade positiva a $(-\infty,0]$. Segundo, as variâncias diferem: a exponencial de taxa
$\lambda$ tem variância $1/\lambda^2$, e $\Var(X\mid Y=0)=1$, enquanto $\Var(X\mid Y=1)=1/9$.

Os dois itens juntos mostram que as hipóteses do LDA são suficientes para o modelo logístico,
mas não necessárias. A regressão logística supõe apenas que o \emph{log-odds} seja afim, e
aqui ele é porque $\log f_c(x)=\log\lambda_c-\lambda_cx$ é afim em $x$ nas duas classes, sem
nada de normal. Para Renata, a diferença é concreta. Com uma amostra grande, o LDA estimaria
$\mu_0=1$, $\mu_1=1/3$ e a variância comum $\tfrac12\big(1+\tfrac19\big)=\tfrac59$, chegaria ao
\emph{log-odds} $0{,}8-1{,}2\,x$ e classificaria como risco alto quem tem $x<2/3$; a fronteira de
Bayes, que a regressão logística estima corretamente, é $x<(\log3)/2\approx0{,}549$. Nesta
população a diferença custa pouco, $0{,}311$ de probabilidade de erro contra $0{,}308$, mas não
diminui com o tamanho da amostra, porque vem da hipótese errada, e não da variabilidade da
estimação.
\end{solucao}


\begin{criterio}
\textbf{(a)} 0,4 pela identidade $\Prob(Y=1\mid\X=\x)=\sigma\big(u(\x)\big)$, pelo Teorema
de Bayes, com $u$ o \emph{log-odds}; 0,3 pelo cancelamento das constantes de
normalização, que a covariância comum torna iguais; 0,6 pela expansão das formas
quadráticas e pelo cancelamento do termo $\x^\top\bm\Sigma^{-1}\x$, com a simetria de
$\bm\Sigma^{-1}$; 0,4 pelos coeficientes, $\bbeta=\bm\Sigma^{-1}(\bm\mu_1-\bm\mu_0)$ e
$\beta_0$, e pela conclusão de que o modelo logístico vale exatamente. Erros de sinal ou
do fator $\tfrac12$ em $\beta_0$ descontam 0,1 cada.

\textbf{(b)} 0,5 pela conta do \emph{log-odds}, $\log3-2x$, e por
$\Prob(Y=1\mid X=x)=\sigma(\log3-2x)$; 0,3 por concluir que o modelo logístico vale, com
$\beta_0=\log3$ e $\beta_1=-2$; 0,5 pelas hipóteses do LDA, que não valem: 0,25 pela falta
de normalidade (o suporte é $(0,\infty)$) e 0,25 pelas variâncias diferentes, $1$ e
$1/9$. A leitura de que as hipóteses do LDA são suficientes mas não necessárias, e a
comparação com o LDA ajustado, não são exigidas.
\end{criterio}

%==============================================================================
\questao{}{3,0 pontos}
%==============================================================================
\pontositens{0,8;1,4;0,8}
% Aula 10 --- o voto de k vizinhos: a variância cai com k, e o 1-NN erra no máximo o dobro do classificador de Bayes.
Otávio é agrônomo e usa o KNN para classificar os talhões de uma fazenda em atacados por uma
praga ($Y=1$) ou não ($Y=0$), a partir de medidas de solo e de clima. Num talhão novo, com
covariáveis $\x$, a probabilidade de praga é $\eta(\x)=\Prob(Y=1\mid\X=\x)=0{,}6$. O KNN
estima essa probabilidade pela proporção de praga entre os $k$ talhões de treino mais
próximos de $\x$, cujos índices formam o conjunto $\mathcal{N}_{\x}^{(k)}$,
\[
  \widehat P=\frac1k\sum_{i\in\mathcal{N}_{\x}^{(k)}}Y_i ,
\]
e declara praga quando $\widehat P>1/2$. Tome $k$ ímpar, para que não haja empate. Para
isolar o papel de $k$, suponha que os rótulos dos vizinhos sejam independentes entre si e
do rótulo $Y$ do talhão novo, e que cada um seja igual a $1$ com probabilidade $0{,}6$, como
se todos os vizinhos estivessem no próprio $\x$.
\begin{itens}
  \item \pts{1,1} Mostre que $\E[\widehat P]=0{,}6$ e $\Var(\widehat P)=0{,}24/k$.
  \item \pts{1,8} Seja $q_k=\Prob(\widehat P>1/2)$. Mostre que a probabilidade de o KNN errar
        o talhão novo é $0{,}6-0{,}2\,q_k$. Calcule $q_5$ e, usando $q_1=0{,}600$,
        $q_{15}\approx0{,}787$ e $q_{51}\approx0{,}926$, a probabilidade de erro para
        $k=1,5,15$ e $51$. Use a desigualdade de Chebyshev para mostrar que $q_k\to1$ e que a
        probabilidade de erro tende a $0{,}4$. Que classificador erra com probabilidade
        $0{,}4$ nesse talhão?
  \item \pts{1,1} Mostre que, com $k=1$ e uma probabilidade $\eta$ qualquer no lugar de
        $0{,}6$, a probabilidade de erro é $2\eta(1-\eta)$, e que
        $\min(\eta,1-\eta)\le2\eta(1-\eta)\le2\min(\eta,1-\eta)$. Interprete. O que acontece
        no outro extremo, com $k=n$?
\end{itens}

\begin{solucao}
Na idealização do enunciado, os rótulos $Y_i$, $i\in\mathcal{N}_{\x}^{(k)}$, são independentes e
têm distribuição de Bernoulli de parâmetro $0{,}6$, e são independentes do rótulo $Y$ do talhão
novo, que também é Bernoulli$(0{,}6)$, porque $\eta(\x)=0{,}6$. O KNN declara praga quando
$\widehat P>1/2$, isto é, quando mais da metade dos $k$ vizinhos tem praga; como $k$ é ímpar,
$k\widehat P$ é um inteiro diferente de $k/2$, e nunca há empate.

\smallskip
\textbf{(a)} O conjunto $\mathcal{N}_{\x}^{(k)}$ tem $k$ elementos, e cada $Y_i$ tem média
$0{,}6$ e variância $0{,}6\cdot0{,}4=0{,}24$. Pela linearidade da esperança,
$\E[\widehat P]=\frac1k\cdot k\cdot0{,}6=0{,}6$. Pela independência, a variância da soma é a
soma das variâncias, e
\[
  \Var(\widehat P)=\frac1{k^2}\sum_{i\in\mathcal{N}_{\x}^{(k)}}\Var(Y_i)
  =\frac{k\cdot0{,}24}{k^2}=\frac{0{,}24}{k} .
\]
Nesta idealização, a proporção de vizinhos com praga é um estimador não viesado de $\eta(\x)$,
e o seu desvio-padrão, $\sqrt{0{,}24/k}$, cai de $0{,}49$ com $k=1$ para $0{,}22$ com $k=5$ e
$0{,}07$ com $k=51$.

\smallskip
\textbf{(b)} O KNN erra se declara praga num talhão sem praga ou se deixa de declarar num
talhão com praga. A decisão $\ghat(\x)=\1{\widehat P>1/2}$ é função só dos rótulos dos
vizinhos, que são independentes de $Y$; por isso
\begin{align*}
  \Prob\big(\ghat(\x)\ne Y\big)
  &=\Prob\big(\ghat(\x)=1\big)\,\Prob(Y=0)+\Prob\big(\ghat(\x)=0\big)\,\Prob(Y=1)\\
  &=0{,}4\,q_k+0{,}6\,(1-q_k)=0{,}6-0{,}2\,q_k .
\end{align*}
Com $k=5$, o número de vizinhos com praga, $5\widehat P$, tem distribuição binomial de
parâmetros $5$ e $0{,}6$, e o KNN declara praga quando ele é pelo menos $3$:
\[
  q_5=\binom53(0{,}6)^3(0{,}4)^2+\binom54(0{,}6)^4(0{,}4)+(0{,}6)^5
     =0{,}3456+0{,}2592+0{,}07776\approx0{,}683 .
\]
Com os valores dados para $q_1$, $q_{15}$ e $q_{51}$, as probabilidades de erro são:
\begin{center}\small
\renewcommand{\arraystretch}{1.25}
\begin{tabular}{@{}lcccc@{}}
\toprule
$k$ & $1$ & $5$ & $15$ & $51$ \\
\midrule
$q_k$ & $0{,}600$ & $0{,}683$ & $0{,}787$ & $0{,}926$ \\
$\Prob\big(\ghat(\x)\ne Y\big)=0{,}6-0{,}2\,q_k$ & $0{,}480$ & $0{,}463$ & $0{,}443$ & $0{,}415$ \\
\bottomrule
\end{tabular}
\end{center}
Para o limite, como $k$ é ímpar, o KNN deixa de declarar praga exatamente quando
$\widehat P<1/2$. Se $\widehat P<1/2$, então $0{,}6-\widehat P>0{,}1$; o evento
$\{\widehat P<1/2\}$ está contido em $\{\abs{\widehat P-0{,}6}\ge0{,}1\}$, e, pela desigualdade
de Chebyshev,
\[
  0\ \le\ 1-q_k=\Prob\big(\widehat P<1/2\big)
  \ \le\ \Prob\big(\abs{\widehat P-0{,}6}\ge0{,}1\big)
  \ \le\ \frac{\Var(\widehat P)}{0{,}1^2}=\frac{24}{k}\ \xrightarrow[k\to\infty]{}\ 0 .
\]
Logo $q_k\to1$, e a probabilidade de erro tende a $0{,}6-0{,}2=0{,}4$. A cota de Chebyshev é
grosseira: só diz algo para $k>24$ e, com $k=51$, garante $1-q_{51}\le0{,}47$, quando o valor é
$0{,}074$. Mas basta para o limite.

O limite $0{,}4$ é a probabilidade de erro do classificador de Bayes nesse talhão: como
$\eta(\x)=0{,}6>1/2$, ele declara praga e erra com probabilidade $1-\eta(\x)=0{,}4$. Nenhuma
decisão em $\x$ faz melhor, porque declarar praga erra com probabilidade $0{,}4$, e não
declarar, com $0{,}6$. Nesta idealização não há viés, aumentar $k$ só reduz a variância de
$\widehat P$, e a probabilidade de erro do KNN se aproxima da de Bayes.

\smallskip
\textbf{(c)} Com $k=1$, o KNN repete o rótulo do vizinho mais próximo, que vale $1$ com
probabilidade $\eta$ e é independente de $Y$. Ele erra se os dois rótulos diferem:
\[
  \Prob\big(\ghat(\x)\ne Y\big)=\eta\,(1-\eta)+(1-\eta)\,\eta=2\eta(1-\eta) .
\]
Seja $m=\min(\eta,1-\eta)$, a probabilidade de erro do classificador de Bayes em $\x$. Então
$1-m=\max(\eta,1-\eta)$, $2\eta(1-\eta)=2m(1-m)$ e $1/2\le1-m\le1$; multiplicando essas duas
desigualdades por $2m\ge0$, obtém-se
\[
  m\ \le\ 2m(1-m)\ \le\ 2m .
\]
Nesta idealização, o $1$-NN erra pelo menos tanto quanto o classificador de Bayes, o que
nenhum classificador evita, e no máximo o dobro; com $\eta=0{,}6$, são $0{,}48$ contra $0{,}40$.
O excesso vem da variância: o $1$-NN copia um único rótulo, sujeito ao mesmo ruído que o rótulo
a prever.

No outro extremo, com $k=n$, o conjunto $\mathcal{N}_{\x}^{(n)}$ é o treino inteiro, qualquer que
seja $\x$. Então $\widehat P=\frac1n\sum_{i=1}^nY_i$ é a proporção de talhões com praga no
treino, que não depende de $\x$, e o KNN atribui a todo talhão a mesma classe, a majoritária no
treino (com $n$ ímpar, sem empate): é um classificador constante, que ignora as covariáveis. A
idealização deixa de valer aqui, porque os vizinhos distantes não estão no próprio $\x$:
$\widehat P$ estima a proporção de praga na população, $\Prob(Y=1)$, e não $\eta(\x)$, e, como
estimador de $\eta(\x)$, tem viés $\Prob(Y=1)-\eta(\x)$, que não diminui com $n$. Entre os dois
extremos, escolher $k$ é equilibrar a variância, que domina com $k$ pequeno, e o viés, que
cresce à medida que os vizinhos se afastam de $\x$.
\end{solucao}


\begin{criterio}
\textbf{(a)} 0,3 pela esperança, pela linearidade; 0,5 pela variância: 0,3 por somar as
variâncias, usando a independência, e 0,2 pela conta, $0{,}24/k$.

\textbf{(b)} 0,4 pela probabilidade de erro $0{,}6-0{,}2\,q_k$, com a independência entre
a decisão e $Y$; 0,3 por $q_5\approx0{,}683$, pela binomial; 0,2 pelas quatro
probabilidades de erro; 0,3 pela desigualdade de Chebyshev, com a inclusão
$\{\widehat P<1/2\}\subset\{\abs{\widehat P-0{,}6}\ge0{,}1\}$ e a cota $24/k$, e pelos
limites $q_k\to1$ e erro $\to0{,}4$; 0,2 por identificar o classificador de Bayes, que erra
com probabilidade $0{,}4$. Usar a desigualdade de Chebyshev sem justificar a inclusão dos
eventos vale 0,1 dos 0,3.

\textbf{(c)} 0,2 pela probabilidade de erro $2\eta(1-\eta)$ com $k=1$; 0,3 pelas
desigualdades $\min(\eta,1-\eta)\le2\eta(1-\eta)\le2\min(\eta,1-\eta)$; 0,1 pela
interpretação: o $1$-NN erra pelo menos tanto quanto o classificador de Bayes e no máximo
o dobro; 0,2 pelo caso $k=n$: o KNN vira o classificador constante, que atribui a classe
majoritária do treino a todo talhão. Dizer que, com $k=n$, a proporção estima $\Prob(Y=1)$,
e não $\eta(\x)$, é bem-vindo, mas não é exigido.
\end{criterio}

\vfill
\begin{center}\small\itshape Boa prova!\end{center}

\end{document}
